% ============================================================================
%  Chapitre 7 — Conclusion et perspectives
%  Cible : 3 à 4 pages. À rédiger J12.
% ============================================================================
\chapter{Conclusion et perspectives}
\label{chap:conclusion}

\section{Rappel de la problématique et de la démarche}
\label{sec:concl-rappel}

\aecrire{Un paragraphe. Reformuler la question de recherche et résumer la
démarche suivie. Ne pas introduire d'élément nouveau.}

\section{Principaux résultats}
\label{sec:concl-resultats}

\aecrire{Reprendre les six contributions annoncées en
section~\ref{sec:intro-contributions} et, pour chacune, énoncer le résultat
obtenu en une à deux phrases, avec les valeurs chiffrées du chapitre~6. La
symétrie entre l'introduction et la conclusion est ce que le jury vérifie en
premier.}

\section{Apports pour l'entreprise}
\label{sec:concl-apports}

\aecrire{Distinguer l'apport technique de l'apport organisationnel : au-delà
de l'outillage, la chaîne modifie les pratiques de l'équipe, rend visible la
qualité et fournit les éléments de preuve attendus par le cadre normatif.}

\section{Apports personnels}
\label{sec:concl-apports-perso}

\aecrire{Une demi-page, factuelle et sans emphase : compétences acquises,
difficultés surmontées, posture d'ingénieur au sein d'une organisation
industrielle.}

\section{Perspectives}
\label{sec:concl-perspectives}

\aecrire{Formuler des perspectives précises et réalistes, issues des limites
du chapitre~6. Éviter les généralités.}

\subsection{Court terme}
\begin{itemize}[nosep]
  \item Application d'une politique de blocage sur les défauts d'analyse
        statique nouvellement introduits.
  \item Exécution de la chaîne à chaque proposition de modification et non
        uniquement selon une planification quotidienne.
  \item Suppression de l'exécution en superutilisateur dans le conteneur de
        compilation.
\end{itemize}

\subsection{Moyen terme}
\begin{itemize}[nosep]
  \item Intégration du banc de test matériel dans la chaîne, afin d'exécuter
        les tests sur cible réelle et non uniquement en compilation native.
  \item Génération et publication d'une nomenclature logicielle normalisée
        associée à chaque livrable.
  \item Signature et attestation de provenance des artefacts produits.
\end{itemize}

\subsection{Long terme}
\begin{itemize}[nosep]
  \item Extension de la chaîne aux autres produits de la gamme, par
        factorisation des workflows réutilisables.
  \item Automatisation de la constitution du dossier de preuves attendu par le
        cadre normatif.
\end{itemize}
